\section{Method: DTA-JEPA}\label{sec:model}

\subsection{Problem setup}

We consider trajectories of observations $o_t\in\R^{n_o}$ (RGB frames plus proprioception),
actions $a_t\in\R^{n_a}$, and a goal observation $o_g$. A latent world model consists of an
encoder $E_s$, an action encoder $E_a$ and a predictor $F$:
\begin{equation}
z_t = E_s(o_t,p_t),\qquad u_t = E_a(a_t),\qquad
(\yhat_{t+1},\Sigma_{t+1}) = F(z_{t-K+1:t},u_{t-K+1:t},u_t),
\end{equation}
trained on reward-free transitions $\mathcal D_{\mathrm{off}}=\{(o_t,a_t,o_{t+1})\}$ by latent
prediction. Goal-conditioned planning minimises a latent cost over an action sequence of
horizon $H$, MPC executes the first $c$ actions, and the loop repeats.

The adaptation problem: after each executed chunk we observe a new transition
$(o_t,a_t,o_{t+1})$ and may update a subset $\Omega$ of parameters before replanning. AdaJEPA
takes exactly one gradient step on $\Omega=\{\text{predictor last block},\text{encoder head}\}$
at a fixed learning rate. DTA-JEPA makes three changes: \emph{what} is updated (a decoupled
fast parameter set), \emph{how much} (an uncertainty-allocated budget), and \emph{for how
long} (state carried across episodes).

\subsection{Encoders and goal conditioning}\label{sec:enc}

The state encoder is a from-scratch convolutional trunk followed by an adaptable projection
head,
\begin{equation}
z_t = \mathrm{LN}\!\Big(W_h\big[\,\mathrm{body}(o_t);\,\phi(p_t)\,\big]\Big),\qquad
\mathrm{body}:\ 64{\times}64{\times}3 \xrightarrow{\ \text{4 conv stages}\ } 128{\times}4{\times}4,
\end{equation}
where $\phi$ is a small MLP over proprioception. The final LayerNorm is not cosmetic: without
it, the cheapest way for $W_h$ to reduce prediction error is to shrink every latent towards
zero, which is exactly the collapse we observed and had to fix at the source
(\S\ref{sec:collapse}). Normalising the output bounds the per-sample scale, so the encoder must
encode \emph{direction}.

Two design choices follow the human-inspired architecture of Cerebro \cite{cerebro-internal}
and LeWM's stability argument \cite{maes-2026-lewm}: (i) no pre-trained semantic encoder and
no EMA---we train everything end-to-end from pixels; (ii) fusion happens \emph{after}
prediction: the goal is encoded with the same encoder, and the interaction between the current
prediction and the goal takes place in the cost function, not by concatenating modalities
before the predictor. All modalities share one latent width $d$.

Proprioception deliberately excludes the manipulated object's pose: including it would let the
model bypass vision entirely and would make the visual-shift suite vacuous. The encoder head
$W_h$ and $\phi$ belong to the fast parameter set; the trunk belongs to the slow set.

\subsection{Dual-timescale recursive predictor}\label{sec:predictor}

\paragraph{Slow module (context).} A short history of $K=3$ latent--action pairs is embedded
into $K$ tokens, $x_i = W_s[z_i;u_i]$, and passed through $L_S$ pre-norm transformer blocks;
the last token is the dynamics context
\begin{equation}
c_t = \mathrm{Transformer}_{L_S}(x_{1:K})_K .
\end{equation}
$c_t$ is the slow state: it summarises "what the dynamics look like here" and is the only
channel through which the slow timescale influences prediction. It is deliberately
\emph{never} updated online.

\paragraph{Fast module (recursive refinement).} The next latent is obtained by recursing in
latent space from the context,
\begin{equation}
\yhat^{(0)} = c_t + u_t,\qquad
\yhat^{(k)} = \yhat^{(k-1)} + g_{\f}\!\big([\yhat^{(k-1)},c_t] + u_t\big),\quad k = 1,\dots,K_t,
\label{eq:refine}
\end{equation}
where $g_{\f}$ is a single pre-norm transformer block, and the update is residual. The fast
parameter set is not fixed by the architecture: our default folds any trained low-rank
adapters into the base weights and adapts the last block directly (which is what our
AdaJEPA baseline does, so the two methods differ \emph{only} in the adaptation policy);
adapting through rank-4 LoRA instead is reported as an ablation (\S\ref{sec:results}). Each depth $k$ carries its own mean head $W_\mu$ and log-variance head $W_{\log\sigma^2}$:
\begin{equation}
\hat\mu^{(k)} = W_\mu \mathrm{LN}\big(g_{\f}^{(k)}\big),\qquad
\log\hat\sigma^{2(k)} = W_{\log\sigma^2}\big(\hat\mu^{(k)}\big).
\end{equation}
The number of refinement steps $K_t\in[K_{\min},K_{\max}]$ is chosen by the controller
(\S\ref{sec:ctrl}); the module also reports a per-sample halting depth $k^\star$, the first
depth at which the relative change $\norm{\yhat^{(k)}-\yhat^{(k-1)}}/\norm{\yhat^{(k-1)}}$
falls below $\epsilon$.

This is the HRM/TRM structure transplanted into a world model: a slow module supplies a
context (HRM's high-level module) and a fast module iterates on a latent answer (TRM's
recursive improvement), with supervision at every depth. Two properties matter for
adaptation: the refinement is parameter-efficient (the fast set is $\sim$0.3M parameters), and
additional \emph{compute} is a first-class axis of adaptation---under shift we can spend more
refinement steps instead of larger gradient steps.

\paragraph{Multi-step rollout and inverse dynamics.} Planning rolls the same module forward,
re-inserting each predicted latent into the history. An auxiliary inverse-dynamics head
$G_\omega([z_t;\sgt(z_{t+1})])\to a_t$ grounds the latent in actions.

\subsection{Uncertainty decomposition and surprise}\label{sec:unc}

We distinguish three quantities, each with a distinct role.

\paragraph{Aleatoric uncertainty} is the predicted diagonal Gaussian:
$\mathcal L_{\mathrm{unc}} = \tfrac12\big[(\yhat-\sgt(z))^\top\Sigma^{-1}(\yhat-\sgt(z)) +
\log\det\Sigma\big]$.

\paragraph{Epistemic uncertainty} is the MC-dropout spread of the R-module,
$\Sigma_{\mathrm{epi}} = \mathrm{Var}_m[\hat\mu_m]$ over $m=4$ stochastic passes, computed only
when the controller requests it.

\paragraph{Surprise} is the Mahalanobis residual of the transition actually observed,
\begin{equation}
s_t = (\yhat_{t+1}-\sgt(z_{t+1}))^{\top}\big(\Sigma_{\mathrm{al}}+\Sigma_{\mathrm{epi}}\big)^{-1}
(\yhat_{t+1}-\sgt(z_{t+1})),
\label{eq:surprise}
\end{equation}
i.e.\ "how implausible is what just happened under the current model". Following LeWM
\cite{maes-2026-lewm}, surprise is the natural detector of events the model cannot represent;
we promote it from a diagnostic to the control signal of the whole adaptation loop.

\subsection{Uncertainty-allocated adaptation controller}\label{sec:ctrl}

Let $\bar s_t$ be an exponential moving average of surprise and let
$r_t = s_t/\bar s_t$ be the relative surprise. Define the allocation
\begin{equation}
\pi_t = \sigma\big(\gamma (r_t-1)\big)\in(0,1),\qquad
U_t = \lfloor U_{\min}+\pi_t(U_{\max}-U_{\min})\rceil,\quad
\eta_t = (1-\pi_t)\eta_{\min}+\pi_t \eta_{\max},
\end{equation}
\begin{equation}
K_t = \lfloor K_{\min}+\pi_t(K_{\max}-K_{\min})\rceil .
\end{equation}
With defaults $[U_{\min},U_{\max}]=[1,3]$, $[\eta_{\min},\eta_{\max}]=[0.2,5]\times\eta_f$ and
$[K_{\min},K_{\max}]=[1,4]$, this yields a concrete statement: \emph{when the model already
predicts the environment, DTA-JEPA is as cheap as AdaJEPA (one step, short recursion); when
the model is wrong, it spends steps, learning rate and refinement depth in proportion to how
wrong it is.} The same $\pi_t$ scales the planning penalties $\beta,\gamma$
(\S\ref{sec:safe}) and the memory write priority (\S\ref{sec:mem}).

\subsection{Dual memory}\label{sec:mem}

\paragraph{Episodic memory} stores latent transitions $(z,u,z')$---never pixels---with
priority
\begin{equation}
\mathrm{prio}(z,u,z') = \lambda_1 s + \lambda_2\,\mathrm{tr}\,\Sigma_{\mathrm{epi}}
+ \lambda_3\,\mathrm{novelty}(z),\qquad
\mathrm{novelty}(z)=\exp\!\big(-{\textstyle\min_m}\norm{z-z_m}^2/\tau\big).
\end{equation}
It is a bounded priority queue, so the model keeps the transitions that were hardest to
predict rather than merely the most recent ones. Adaptation retrieves the top-$k$ memories by
a Gaussian similarity kernel and adds their prediction loss:
\begin{equation}
\mathcal L_{\mathrm{on}} = \mathcal L_{\mathrm{ada}}(\mathcal B_t)
+ \lambda_{\mathrm{ret}}\,\mathcal L_{\mathrm{ada}}(\mathcal M_{\mathrm{ret}}).
\end{equation}

\paragraph{Parametric memory} is a skill bank: fast-parameter snapshots indexed by an
environment fingerprint $\mathrm{fp}$ (the mean latent of the episode's first frames). At
episode start the nearest snapshot warm-starts $\f$. This is what turns per-episode adaptation
into cross-episode learning: returning to a previously seen environment is no longer a
cold start.

\paragraph{Parameterisation is an experimental control, not a contribution.} Because
AdaJEPA's own ablations show that rank-4 LoRA adapters recover less of the frozen-to-adapted
gap than direct updates of the same block, evaluating a LoRA-based method against a
direct-update baseline would confound the adaptation \emph{policy} with the adaptation
\emph{parameterisation}. We therefore fold the trained adapters into the base weights
(numerically exact: the folded and unfolded models agree to $3\times10^{-8}$ in prediction
loss) and give both methods the identical fast set---the predictor's last block plus the
encoder head, with the encoder trunk and the S-module frozen---so that measured differences
are attributable to allocation, memory, safety and refinement depth alone. The LoRA
parameterisation is then a separate row in the ablation table, and it costs accuracy there,
which is consistent with AdaJEPA's finding.

\paragraph{Slow consolidation} happens once per episode, on the episodic memory, with an EWC
penalty: $\ts \leftarrow \ts - \eta_s\nabla_{\ts}[\mathcal L_{\mathrm{ada}}(\mathcal M_e) +
\lambda_{\mathrm{EWC}}\sum_i F_i(\theta_{s,i}-\theta_{s,i}^\star)^2]$ with $\eta_s = 10^{-6}$. The fast
parameters are never consolidated into the slow ones during an episode, and the slow ones are
never touched by online gradients.

\subsection{Safety layer}\label{sec:safe}

Four mechanisms bound the damage of a bad update.

\begin{enumerate}[leftmargin=1.2em,itemsep=2pt]
\item \textbf{Anchor ball.} $\norm{\f-\f^0}_2\le\epsilon_f$, enforced by projection after every
step: $\f \leftarrow \f^0 + \frac{\epsilon_f}{\norm{\f-\f^0}}(\f-\f^0)$.
\item \textbf{Validation-gated rollback.} The oldest entries of the replay buffer form a
held-out slice; if $\mathcal L_{\mathrm{val}}^{\mathrm{new}} >
\mathcal L_{\mathrm{val}}^{\mathrm{old}}+\varepsilon$, the update is reverted and the event is
logged.
\item \textbf{Immutable slow anchor.} $\ts$ (encoder trunk and S-module) receives no online
gradient: the pretrained dynamics context cannot be corrupted by a bad episode.
\item \textbf{Support penalty in planning.} Let $\tau_{0.95}$ be the 95th percentile of the
training latent Mahalanobis radius. The planner minimises
\begin{equation}
C(a)=\sum_{k=1}^{H}\Big[\underbrace{\norm{\yhat_k-z_g}^2}_{\text{goal}}
+\underbrace{\beta\,\mathrm{tr}\,\Sigma_k}_{\text{uncertainty}}
+\underbrace{\gamma\max\big(0,\ \mathrm{Maha}(\yhat_k)-\tau_{0.95}\big)}_{\text{support}}\Big],
\label{eq:cost}
\end{equation}
so imagined trajectories are penalised for leaving the region of latent space the model has
evidence about. Frozen and AdaJEPA use $\beta=\gamma=0$ (AdaJEPA's eq.~(3)).
\end{enumerate}

\subsection{Training}\label{sec:train}

\paragraph{Stage 1: pretraining.} A two-term objective in the spirit of LeWM, extended with
deep supervision:
\begin{equation}
\mathcal L = \underbrace{\sum_{k=1}^{K_{\max}} w_k \mathcal L_{\mathrm{mse}}\big(\yhat^{(k)},
\sgt(z_{t+1})\big)}_{\text{deep supervision}}
+ \underbrace{\sum_{j=1}^{J}\mathcal L_{\mathrm{mse}}\big(\yhat_{t+j},\sgt(z_{t+j})\big)}_{\text{multi-step rollout}}
+ \lambda_{\mathrm{unc}}\mathcal L_{\mathrm{unc}}
+ \lambda_{\mathrm{reg}}\mathcal R(z)
+ \lambda_{\mathrm{inv}}\mathcal L_{\mathrm{inv}},
\end{equation}
with $J=4$ future steps, increasing depth weights $w_k$ and the regulariser
$\mathcal R(z)=\mathrm{SIGReg}(z)+\mathrm{VarFloor}(z)$: SIGReg is the Epps--Pulley
isotropic-Gaussian test on random 1-D projections used by LeJEPA/LeWM
\cite{balestriero-2025-lejepa}, and $\mathrm{VarFloor}(z)=\frac1d\sum_i\mathrm{relu}(1-\mathrm{std}_i(z))$
is a VICReg-style \cite{bardes-2022-vicreg} hinge that keeps the statistic informative in the
collapse regime. All targets are stop-gradient.

\paragraph{Stage 2: meta-training.} The fast-parameter initialisation $\f^0$ is meta-learned
with first-order MAML (FOMAML) over tasks defined by holding out one shape (PushBlock) or one
layout group (maze):
\begin{equation}
\f' = \f^0-\eta_f\nabla_{\f}\mathcal L_{\mathrm{ada}}(\mathcal S_\tau),\qquad
\f^0 \leftarrow \f^0-\eta_{\mathrm{meta}}\nabla_{\f^0}\mathcal L_{\mathrm{ada}}(\mathcal Q_\tau;\f'),
\end{equation}
so that a single online step lands in a useful region. Slow parameters are frozen during
meta-training.

\paragraph{Stage 3: deployment.} Algorithm~\ref{alg:dtajepa} summarises the online loop.

\begin{figure}[t]
\centering
\fbox{\parbox{0.96\textwidth}{\footnotesize
\textbf{Algorithm 1: DTA-JEPA online closed-loop adaptation}
\begin{tabbing}
\hspace{1.1em}\=\hspace{2.2em}\=\hspace{2.2em}\=\kill
\textbf{input}: $\ts^0$ (immutable anchor), $\f^0$, $\mathcal M_e$, $\mathcal M_p$, $o_g$, hyper-parameters\\
$\f \leftarrow \mathrm{Retrieve}(\mathcal M_p,\mathrm{fp})$;\quad $z_g\leftarrow E_s(o_g)$;\quad $\mathcal B\leftarrow\emptyset$\\
\textbf{for} $t=0,\dots,T-1$ \textbf{do}\\
\> \# 1. plan with the current model\\
\> $c_t \leftarrow \mathrm{S}(z_{t-K+1:t},u_{t-K+1:t})$;\quad
   $\yhat_{1:H} \leftarrow \mathrm{Rollout}(c_t,a_{t:t+H-1},K_t)$\\
\> $a^\star_{t:t+H-1} \leftarrow \arg\min C(a)$ \hfill (eq.~\ref{eq:cost})\\
\> \# 2. execute a chunk, observe a labelled transition\\
\> execute $a^\star_{t:t+c-1}$;\quad observe $o_{t+1}$;\quad
   $\mathcal B \leftarrow \mathcal B\cup\{(o_t,a_t,o_{t+1})\}$ (recent-$N$)\\
\> \# 3. surprise and allocation\\
\> $s_t \leftarrow$ Mahalanobis residual (eq.~\ref{eq:surprise});
   $U_t,\eta_t,K_t \leftarrow$ Controller$(s_t)$\\
\> \# 4. bounded adaptation of $\f$ only\\
\> $\mathcal M_{\mathrm{ret}}\leftarrow \mathrm{TopK}(\mathcal M_e)$;\quad
   \textbf{for} $u=1..U_t$ \textbf{do} $\f \leftarrow \f-\eta_t\nabla_{\f}[\mathcal L_{\mathrm{ada}}(\mathcal B)
   +\lambda_{\mathrm{ret}}\mathcal L_{\mathrm{ada}}(\mathcal M_{\mathrm{ret}})]$ \textbf{end}\\
\> $\f \leftarrow$ project($\f$, anchor ball);\quad \textbf{if} validation worsens \textbf{then} rollback\\
\> $\mathcal M_e.\mathrm{add}(z_t,u_t,z_{t+1};s_t)$\\
\textbf{end}\\
\textbf{episode end} (slow timescale):\\
\> $\mathcal M_p.\mathrm{add}(\mathrm{fp},\f)$;\quad
   $\ts \leftarrow \ts-\eta_s\nabla_{\ts}[\mathcal L_{\mathrm{ada}}(\mathcal M_e)+\lambda_{\mathrm{EWC}}\,
   \text{EWC}(\ts)]$
\end{tabbing}}}
\caption{DTA-JEPA's online loop. Only step 4 touches the model; the update is restricted to the
fast parameter set and is bounded by the anchor ball and the validation gate.}
\label{alg:dtajepa}
\end{figure}

\subsection{Why this should help: two propositions}

\begin{proposition}[Anchor-bounded drift]
Let the slow anchor $\ts^0$ define the reference dynamics. If (i) online updates are confined
to $\f$ with $\norm{\f-\f^0}\le\epsilon_f$, and (ii) the predictor is Lipschitz in $\f$ with
constant $L_P$ on the visited latent support, then for every visited $(z,u)$
$\norm{F_{\ts^0,\f}(z,u)-F_{\ts^0,\f^0}(z,u)}\le L_P\epsilon_f$. In particular, an
in-distribution environment whose prediction error under $\f^0$ is already below
$L_P\epsilon_f$ cannot be made worse without bound by adaptation, whereas an unconstrained
update of the pretrained weights carries no such guarantee.
\end{proposition}

\begin{proposition}[Compute allocation]\label{prop:alloc}
Suppose the refinement operator contracts towards the target at a rate depending on the shift:
$\norm{\yhat^{(k)}-\sgt(z_{t+1})}\le \rho^k\norm{\yhat^{(0)}-\sgt(z_{t+1})}+\delta$, with
$\rho<1$ and approximation error $\delta$. Then the excess prediction error achievable with
depth $K_t$ is monotone decreasing in $K_t$ and the number of steps needed to reach a target
error $\varepsilon$ grows only logarithmically in the initial residual. Consequently, when
surprise is high (large initial residual) increasing $K_t$ is the cheapest available currency,
while when surprise is low ($\rho^k$ residual already below $\varepsilon$) all $K_t>1$ are
wasted compute.
\end{proposition}

Both propositions are statements about our own mechanism, not about the task: they justify the
two design choices (anchor ball on a decoupled fast set; computational allocation governed by
residual size), and they predict the empirical pattern we look for in \S\ref{sec:results}---a
gain that is largest out of distribution and smallest in distribution, at equal or lower
average cost.
